\section{Implementation and reproducibility}
\label{sec:implementation}

This section describes how the methodology is realised in software. It is not
an appendix to the science: it is the part of the paper that makes the
preceding claims \emph{checkable} rather than merely asserted. A methodology
whose values live in prose can drift from the code that applies them, and a
reader has no way to detect the drift. The arrangement described here is
designed so that such drift fails the build.

\subsection{Configuration as data}
\label{sec:config-as-data}

Every methodology value lives in schema-validated, version-stamped YAML under
\texttt{config/}:

\begin{table}[H]
\centering
\caption{The methodology configuration files.}
\label{tab:config-files}
\small
\begin{tabularx}{\textwidth}{@{}l>{\raggedright\arraybackslash}X@{}}
\toprule
\textbf{File} & \textbf{Content} \\
\midrule
\texttt{nbs\_typologies.yaml} &
Both levels of the library (\Cref{sec:archetypes}): the \num{18} cooling
archetypes with their base cooling scores, temperature envelopes, evidence
confidence, suitability conditions, co-benefit defaults, output caveats, and a
\texttt{sources} array carrying for each citation the source key and the
specific finding it supports; and the \num{110} catalogue entries with their
identity, family, kind, primary cooling mechanism, availability, the archetype
each inherits, and the \num{21} per-entry suitability overrides. \\
\addlinespace
\texttt{availability.yaml} &
The gating policy of \Cref{sec:availability}: what each of the four site
conditions means, what an \emph{unanswered} condition means, what each
assessment scale offers, the package size warning threshold, and the land-use
mapping counts. It carries no performance value and feeds no score. \\
\addlinespace
\texttt{weights.yaml} &
Every weighted sum in \Cref{sec:formulas}, with its rationale, plus the
sensitivity-analysis perturbation parameter. \\
\addlinespace
\texttt{adjustment\_factors.yaml} &
The condition-level to factor table, the four derivation rules of
\Cref{tab:derivation}, and the climate suitability matrix of
\Cref{tab:climate-matrix}. \\
\addlinespace
\texttt{input\_mapping.yaml} &
Qualitative to numeric normalisation, the inverted scale, the LST
normalisation, and the heat-index improvement buckets. \\
\addlinespace
\texttt{energy\_model.yaml} &
The cooling-energy derivation of \Cref{sec:energy}: the formula, the
sensitivity interval of \numrange{0.02}{0.04} per \si{\degreeCelsius} with its
citation, the three preconditions, the reported statuses, and the recorded
limitation on transferability. \\
\addlinespace
\texttt{country\_defaults.yaml} &
Grid emission factor source, licence, and per-country schema; energy price
policy. Deliberately contains no cost values. \\
\addlinespace
\texttt{recommendation\_templates.yaml} &
The deterministic text fragments from which recommendations are composed, and
the conditions under which climate-dependent caveats fire. \\
\addlinespace
\texttt{derived\_scores.yaml} &
The derivation rules fixed at version 2026.08.01: suitability and equity
sub-indicator rules, cost-feasibility payback brackets, investment-readiness
rules, time-to-benefit buckets, the confidence field-to-block mapping, and the
score interpretation bands. \\
\bottomrule
\end{tabularx}
\end{table}

The scoring code contains no typology value, no weight, and no threshold.
Changing the methodology never requires changing code, which has three
consequences. A deployment can rebalance the weights it disagrees with ---
including the equity-forward stance of \Cref{sec:effective-weights} --- without
forking the software. A reviewer can read the complete parameterisation of the
tool in nine configuration files rather than by tracing constants through
source code.
And the diff of a methodology change is legible: it shows values moving, not
logic moving.

\subsection{What is implemented at version \methver{}}
\label{sec:implementation-status}

\begin{keynote}
\noindent This section distinguishes throughout between what the repository
\emph{contains} at methodology version \methver{} and what it \emph{specifies}
for the next development phase. The distinction is stated because a reviewer
can and should check it by cloning the repository, and because a paper that
described planned software in the present tense would forfeit precisely the
credibility this section exists to establish.

\textbf{Implemented and running in continuous integration:} the configuration
layer --- schema-validated loading of all nine methodology files, version
lock-step enforcement, and the citation cross-check against the bibliography ---
the six evidence rules of \Cref{sec:evidence-rules}; the calculation engine
itself, comprising the input and result schemas, one scoring module per formula
family, the availability matrix, the package combination rules of
\Cref{sec:packages}, branched confidence, the deterministic recommendation
composer, and the orchestrating entry point, under a \SI{100}{\percent}
line-coverage gate with \num{24} hand-verified golden scenarios, a
byte-identical determinism test, and the executed sensitivity analysis of
\Cref{sec:sensitivity-results}; the report and workbook export; the HTTP API
with its explicit stored-draft migration (\Cref{sec:governance}); and the web
interface.

The claims in this section are about what the repository contains, and a
reviewer should check them by cloning it. Where this paper describes a
behaviour the engine is \emph{specified} to have rather than one it exhibits,
it says so.

Every number in the worked example of \Cref{sec:worked-example} was first
computed by hand from the specification and is now also asserted, digit for
digit, by the engine's golden-scenario suite --- the hand derivation is the
expected value, never the engine's own output.
\end{keynote}

\subsection{The engine contract}

The calculation engine is specified as a pure function

\[
\texttt{run\_assessment}(\textit{AssessmentInput},\ \textit{MethodologyConfig})
\;\longrightarrow\; \textit{AssessmentResult}
\]

with no input or output, no network access, no global mutable state, and no
stochastic component. Determinism is therefore a structural property of the
design rather than a behaviour that testing hopes to confirm: there is no
source of variation for a test to detect. A determinism test is nonetheless
run, on the principle that a structural guarantee worth relying on is worth
verifying.

Each result records both the engine version and the methodology version that
produced it, and comparisons between results computed under differing
methodology versions raise a warning rather than being silently rendered side
by side.

The contract was published in advance of its implementation, which is the point
at which it could still be criticised cheaply, and it has not been relaxed
since. A reviewer who believes it is too strong, too weak, or unenforceable is
invited to say so.

\subsection{Machine-enforced evidence rules}
\label{sec:evidence-rules}

The project's evidence policy is a build gate, not a review convention.
Configuration loading fails --- and continuous integration fails with it --- if
any of the following holds.

\begin{enumerate}
  \item \textbf{A performance value carries no citation.} The schema requires a
  non-empty \texttt{sources} array on every archetype, and each entry in it must
  carry both a source key and the finding it supports. Because catalogue entries
  inherit their values, the check runs over the \emph{resolved} library, so an
  entry cannot acquire an uncited value through inheritance either.

  \item \textbf{A citation points at an unrecorded source.} Every
  \texttt{sources[].key} appearing anywhere in the configuration tree is
  collected recursively and checked against the set of keys declared in the
  bibliography. A citation naming a source the bibliography does not record is
  a hard error. This is the rule that makes the traceability claim of
  \Cref{sec:commitments} enforceable: a value cannot cite something that does
  not exist.

  \item \textbf{The configuration files fall out of version lock-step.} Every
  file carries a top-level date-stamped \texttt{version}, and all must agree. A
  methodology version moves as a single unit or not at all.

  \item \textbf{A temperature envelope is inverted or implausible.} The maximum
  must not fall below the minimum, and no resolved entry may declare a maximum above
  \qty{3.5}{\degreeCelsius}. The upper guard encodes a methodology judgment: no
  retrieved source supports a daytime pedestrian-level air temperature
  reduction above roughly \qty{3}{\degreeCelsius} for a single site-level
  intervention, so a larger value in configuration indicates a calibration
  error rather than new evidence. Raising the guard requires changing the test,
  which makes the change visible in review.

  \item \textbf{The low-confidence declarations are removed.} A test asserts
  that the green fa\c{c}ade / living wall, bioretention, the enclosed courtyard,
  the vegetated shade structure, and both productive archetypes remain declared
  as low evidence confidence, guarding against a future edit quietly promoting a
  weakly evidenced class.

  \item \textbf{Default costs appear.} A test asserts that the country defaults
  ship no cost or emission values pending verification, and that the typology
  library carries no per-square-metre unit cost. The cost policy of
  \Cref{sec:costs} holds in configuration, not only in prose.
\end{enumerate}

The practical effect is that the integrity claims made in this paper degrade
loudly rather than quietly. A contributor who adds a typology without a
citation, or who bumps one configuration file without the others, does not
produce a subtly wrong tool; they produce a red build.

\subsection{Testing and continuous integration}
\label{sec:testing}

\paragraph{The configuration layer.} The test suite asserts that all \num{18}
archetypes and all \num{110} entries load, resolve, and validate, that every
entry names an archetype that exists, that every configuration file shares one
version stamp, that every resolved entry carries at least one citation with a
stated finding, that every source key resolves against the bibliography, that
temperature envelopes are ordered and below the \qty{3.5}{\degreeCelsius}
guard, that the six low-confidence archetypes remain declared as such, and that
no default cost or unverified emission value ships. Negative cases are tested
explicitly: an invented citation key, a version mismatch, an uncited class, and
an inverted envelope must each fail the load.

These are the six evidence rules of \Cref{sec:evidence-rules}. Continuous
integration runs linting, format checking, strict static type checking, and the
full test suite --- configuration, scoring, confidence, recommendation,
scenarios, determinism --- with the \SI{100}{\percent} engine coverage gate on
every push and pull request, across supported Python versions.

\paragraph{The engine.} Unit tests cover every scoring module, with the
\SI{100}{\percent} engine line-coverage target enforced as a build gate;
\num{24} golden scenarios, each a stored input with a hand-verified expected
output and its derivation, serve as regression armour against unintended
methodology drift and as the basis of the sensitivity analysis of
\Cref{sec:sensitivity-results}; and a determinism test asserts byte-identical
results across repeated runs and freshly loaded configurations.

\paragraph{Three methodology properties are asserted as tests rather than
stated as prose}, because each is a rule a plausible future edit would break.
\emph{Availability feeds no score:} one site is scored twice, with all four
gating questions answered and with none of them answered, and the two results
must be byte-identical (\Cref{sec:availability}). \emph{The two water
archetypes are never interpolated by area:} the same constructed wetland scored
on two sites differing in area by a factor of \num{2500} must report an
identical ceiling (\Cref{sec:small-water}). And \emph{the area-scale shading
figure does not transfer to pergola scale:} the vegetated shade structure must
not cite it, while the permeable shaded hardscape must
(\Cref{sec:shade-structure}). The nine published availability situation counts and the four
land-use totals --- including the \num{72} entries that name \emph{school} in
their land-use list --- are likewise asserted, so a land use silently losing
its interventions fails the build.

The ordering is deliberate rather than incidental. The evidence rules were
implemented before the scoring code they protect, so that no scoring value can
enter the repository uncited at any point in the project's history --- not even
temporarily, during development.

\subsection{Methodology versioning stamped into every result}

The methodology version --- \methver{} at the time of writing --- stamps this
document, the Methodology Report in the repository, and every configuration
file, and is recorded, together with the engine version, in every assessment
result the engine produces. Lock-step across the configuration files is
enforced by the build.

This has a consequence that matters for how the tool behaves over time.
Existing assessments are \emph{never} to be silently recomputed. A result
retains the version that produced it, and the interface indicates when a newer
methodology version is available, leaving the decision to re-run with the user.
A city
whose priority list was produced under version \methver{} can therefore explain,
years later, exactly which rules generated it --- because those rules are a
tagged, public, immutable artefact rather than the current state of a codebase.

\subsection{What this arrangement does and does not guarantee}

It is worth being precise about the limits of these guarantees, since this
section makes the strongest reproducibility claims in the paper.

The arrangement guarantees that every methodology value in the repository
carries a citation to a recorded source, that no citation names a source the
bibliography does not hold, that a change to any value is versioned and visible
across every file at once, and --- since the engine and its coverage gate exist
--- that the tool applies exactly the methodology described in this document.
These are properties of internal consistency and traceability, not of
correctness.

It does \emph{not} guarantee that the methodology is correct, that the adopted
envelopes are the right ones, or that the cited sources support the weight
placed on them. Those are scientific questions, and no build gate can settle
them. What the arrangement does is make them \emph{answerable} by an external
reviewer: the complete parameterisation is public, each value names its source,
and \Cref{app:verification} states what was actually verified about each
source. The engineering makes the science auditable; it does not make it right.
